

\ch{Appendix B: Statistical Analysis and Methods}\label{appendix:statistical-methods}

This appendix documents the complete statistical analysis underlying Chapter~4 and demonstrates each computation functionally. For every procedure it states the method and its originating reference, gives the estimator or test statistic in closed form, and then works a representative case through to its numerical result using the survey data. The detailed tabulations referenced below appear as tables and figures in Chapter~4; the purpose here is to show the arithmetic that produced them, so that every reported value is both traceable to its definition and reproducible by hand.

\section{Reproducibility and Tooling}\label{sec:appb-reproducibility}

A fully scripted Python pipeline (NumPy, pandas, and SciPy, with statsmodels for the regression and differential-item-functioning models and semopy for the confirmatory measurement model) produced all statistics, tables, and figures from the response file. The pipeline computes every quantity reported here and writes the results, figures, and \LaTeX{} table fragments deterministically, so that the analysis is reproducible from the raw data. A companion computational notebook reproduces the same results with the mathematics shown inline. Sample standard deviations and variances use the unbiased $(n-1)$ denominator throughout, and the worked figures below are carried to four decimal places before rounding so that intermediate steps can be checked.

\paragraph{Data and materials availability.} A public reproducibility artifact accompanies this dissertation. The survey instrument, the analysis code, and a synthetic dataset that reproduces the structure of the analyses reported in Chapter~4 are openly available at \url{https://doi.org/10.17605/OSF.IO/VKDTE}. Because the survey was administered anonymously and received an Institutional Review Board exempt determination, the deposited dataset is synthetic and contains no real response record.

\section{Data Preparation and Sample}

The survey platform recorded $1{,}600$ submissions. A pre-registered screen removed $188$ rapid or abandoned cases (completion times under three minutes, open longer than five days, or answering at most the two Section~1 questions). A post-hoc forensic authenticity screen then removed one further submission belonging to a coordinated automated cluster, in which a single identical response vector had been submitted $71$ times in a roughly one-minute burst, leaving an analytic sample of $N=1{,}411$. Decoding mapped Likert items from their anchored string form to the integer scale $\{1,2,3,4,5\}$. Complete data across all twenty-three Likert items (the twenty-two agreement-scale items together with the familiarity item Q1.1) are present for $1{,}390$ respondents ($98.5\%$); item-level missingness never exceeds six cases, and available-case analysis handled the remaining gaps.

\section{Descriptive Statistics}

For each Likert item the sample mean and unbiased sample standard deviation are
\begin{equation}
  \bar{x} = \frac{1}{n}\sum_{i=1}^{n} x_i,
  \qquad
  s = \sqrt{\frac{1}{n-1}\sum_{i=1}^{n}\left(x_i-\bar{x}\right)^2}.
\end{equation}

\paragraph{Worked example (familiarity, Q1.1).} The familiarity item was answered by all $n=1{,}411$ respondents with the frequency distribution $f_1=700$, $f_2=404$, $f_3=59$, $f_4=192$, $f_5=56$. Because the data are discrete, the mean is computed from the frequencies,
\begin{equation}
  \bar{x} = \frac{\sum_{v=1}^{5} v\,f_v}{n}
          = \frac{(1)(700)+(2)(404)+(3)(59)+(4)(192)+(5)(56)}{1411}
          = \frac{2733}{1411} = 1.9369 .
\end{equation}
The corresponding sum of squared deviations gives the variance and standard deviation,
\begin{equation}
  s^2 = \frac{1}{n-1}\sum_{v=1}^{5} f_v\left(v-\bar{x}\right)^2 = 1.4364,
  \qquad s = \sqrt{1.4364} = 1.1985 .
\end{equation}
These match the values reported for Q1.1 in Table~\ref{tab:desc} ($\bar{x}=1.94$, $s=1.20$). Table~\ref{tab:desc} gives the full set of item-level descriptives.

\subsection{Agreement Proportions and the Wilson Interval}

The agreement proportion is $\hat{p}=k/n$, the share of responses at or above $4$. The Wilson score method~\parencite{Wilson_1927} formed each $95\%$ confidence interval for a proportion and remains well-behaved near the response ceiling,
\begin{equation}
  \frac{\hat{p}+\dfrac{z^2}{2n}\;\pm\; z\sqrt{\dfrac{\hat{p}\left(1-\hat{p}\right)}{n}+\dfrac{z^2}{4n^2}}}{1+\dfrac{z^2}{n}},
  \qquad z=z_{0.975}=1.96 .
\end{equation}

\paragraph{Worked example (Q1.1 agreement).} For Q1.1, $k=248$ respondents answered $4$ or $5$, so $\hat{p}=248/1411=0.1758$ with $n=1411$ and $z=1.96$. The denominator and the two half-terms are
\begin{equation}
  1+\frac{z^2}{n}=1.0027,\qquad
  \hat{p}+\frac{z^2}{2n}=0.1771,\qquad
  z\sqrt{\frac{\hat{p}(1-\hat{p})}{n}+\frac{z^2}{4n^2}}=0.0199 .
\end{equation}
Combining gives the interval
\begin{equation}
  \frac{0.1771 \pm 0.0199}{1.0027} = [\,0.1568,\; 0.1965\,] = [\,15.7\%,\; 19.6\%\,],
\end{equation}
which matches the Wilson bounds reported for Q1.1 in Table~\ref{tab:desc}.

\section{Composite Scores and Internal Consistency}

Three composites were formed by averaging their constituent items: Awareness (Q1.1, Q2.1--Q2.6), IT~Value (Q3.1, Q3.3--Q3.5, Q4.1--Q4.4), and Security~Value (Q3.2, Q3.6, Q5.1--Q5.6). Internal-consistency reliability was assessed with Cronbach's $\alpha$~\parencite{Cronbach_1951}, where $\sigma^2_{j}$ is the variance of item $j$ among $k$ items and $\sigma^2_{T}$ the variance of the summed score,
\begin{equation}
  \alpha = \frac{k}{k-1}\left(1-\frac{\sum_{j=1}^{k}\sigma^2_{j}}{\sigma^2_{T}}\right).
\end{equation}

\paragraph{Worked example (Awareness composite, $k=7$).} The seven item variances are
\begin{equation}
\begin{aligned}
  &\sigma^2_{Q1.1}=1.4364,\;\; \sigma^2_{Q2.1}=0.3295,\;\; \sigma^2_{Q2.2}=0.3192,\;\; \sigma^2_{Q2.3}=0.3215,\\
  &\sigma^2_{Q2.4}=0.3352,\;\; \sigma^2_{Q2.5}=0.5155,\;\; \sigma^2_{Q2.6}=0.3295,
\end{aligned}
\end{equation}
which sum to $\sum_j \sigma^2_j = 3.5869$. The variance of the seven-item summed score over complete cases is $\sigma^2_{T}=11.1691$. Substituting,
\begin{equation}
  \alpha = \frac{7}{6}\left(1-\frac{3.5869}{11.1691}\right)
         = (1.1667)(1-0.3211) = (1.1667)(0.6789) = 0.792 .
\end{equation}
The IT~Value ($\alpha=0.858$) and Security~Value ($\alpha=0.858$) composites are computed identically; all three exceed the conventional $\alpha\ge0.70$ threshold (Table~\ref{tab:comp}). The $95\%$ confidence interval for each composite mean follows from Student's $t$ distribution, $\bar{x}\pm t_{0.975,\,n-1}\,s/\sqrt{n}$.

\section{Tests of Normality}

The Shapiro--Wilk test assessed each composite for normality~\parencite{Shapiro_Wilk_1965}. The statistic compares the squared correlation between the ordered sample and the expected normal order statistics,
\begin{equation}
  W = \frac{\left(\sum_{i=1}^{n} a_i x_{(i)}\right)^2}{\sum_{i=1}^{n}\left(x_i-\bar{x}\right)^2},
\end{equation}
where $x_{(i)}$ are the order statistics and $a_i$ the tabulated normal weights. Every composite departed from normality (Awareness $W=0.978$; IT~Value $W=0.957$; Security~Value $W=0.965$; each $p<.001$). Because the test is highly powered at this sample size and the composites are left-skewed, the analysis adopted non-parametric procedures for inference and foregrounded effect sizes over $p$-values.

\section{Effect Sizes}

Cohen's $d$ standardizes the mean difference by the pooled standard deviation~\parencite{Cohen_1988},
\begin{equation}
  d = \frac{\bar{x}_1-\bar{x}_2}{s_p},
  \qquad
  s_p = \sqrt{\frac{(n_1-1)s_1^2+(n_2-1)s_2^2}{n_1+n_2-2}}.
\end{equation}

\paragraph{Worked example (Awareness, IT vs IA).} The Information Technology group ($n_1=501$) has mean $\bar{x}_1=4.379$ and $s_1=0.38$; the Information Assurance group ($n_2=501$) has $\bar{x}_2=3.732$ and $s_2=0.41$. The pooled standard deviation is
\begin{equation}
  s_p = \sqrt{\frac{(500)(0.38^2)+(500)(0.41^2)}{501+501-2}}
      = \sqrt{\frac{72.200+84.050}{1000}} = \sqrt{0.15625} = 0.3953,
\end{equation}
so that
\begin{equation}
  d = \frac{4.379-3.732}{0.3953} = \frac{0.647}{0.3953} = 1.64 .
\end{equation}
This is the value reported for Awareness in Table~\ref{tab:field}. A $95\%$ confidence interval for $d$ was obtained from its large-sample standard error~\parencite{Hedges_Olkin_1985},
\begin{equation}
  \mathrm{SE}_d = \sqrt{\frac{n_1+n_2}{n_1 n_2}+\frac{d^2}{2\,(n_1+n_2)}}
              = \sqrt{\frac{1002}{251001}+\frac{1.64^2}{2(1002)}} = 0.0730,
\end{equation}
so that $d \pm 1.96\,\mathrm{SE}_d = 1.64 \pm 0.14 = [\,1.50,\;1.78\,]$, the interval reported in Table~\ref{tab:field}. The IT~Value and Security~Value intervals ($[1.79,2.09]$ and $[1.79,2.09]$) were computed identically; all three lie well above the conventional ``large'' threshold of $0.8$, so the field difference is both very large and precisely estimated. The rank-biserial correlation, derived from the Mann--Whitney statistic, expresses the same comparison on the ordinal scale~\parencite{Kerby_2014}. The statistic entering the formula is $U_1$ for the Information Technology group, the reference group, computed from that group's rank sum; under the convention $r_{rb}=1-2U_1/(n_1 n_2)$, a negative value indicates that the reference group ranks higher. Because Information Technology respondents score above Information Assurance respondents, $U_1 = 218{,}766$ exceeds $n_1 n_2/2 = 125{,}500.5$ and the coefficient comes out negative,
\begin{equation}
  r_{rb} = 1-\frac{2U_1}{n_1 n_2}
         = 1-\frac{2(218766)}{(501)(501)}
         = 1-\frac{437532}{251001} = -0.743 .
\end{equation}
The sign therefore encodes direction (Information Technology higher), consistent with the positive Cohen's $d$ above and with the sign convention stated in Chapter~4, while the magnitude $0.743$ carries the effect size.

\section{Group Comparison by Professional Field}

The Mann--Whitney $U$ test~\parencite{Mann_Whitney_1947} compared Information Technology ($n=501$) and Information Assurance/Cyber Security ($n=501$) respondents on each composite; the test evaluates whether one group tends to rank above the other without assuming normality. With $R_1$ the rank sum of the first group in the pooled ranking,
\begin{equation}
  U_1 = R_1 - \frac{n_1(n_1+1)}{2}, \qquad U_2 = n_1 n_2 - U_1 .
\end{equation}
The significance computation may take the smaller of $U_1$ and $U_2$; the rank-biserial effect size worked in the preceding section instead uses $U_1$ for the Information Technology group directly, so that its sign encodes which group ranks higher. Welch's $t$ was reported alongside for continuity with the parametric literature. Differences are large on all three composites (Table~\ref{tab:field}): Awareness $d=1.64$ ($r_{rb}=-0.743$, worked above), IT~Value $d=1.94$, and Security~Value $d=1.94$, with Information Technology exceeding Information Assurance throughout. At the item level within the Awareness family, the Holm--Bonferroni step-down correction controlled the family-wise error rate~\parencite{Holm_1979}: the $m$ ordered $p$-values $p_{(1)}\le\cdots\le p_{(m)}$ are compared in turn against $\alpha/(m-i+1)$, and the procedure stops at the first $i$ for which the inequality fails, retaining all remaining hypotheses. Because the unadjusted $p$-values are already extreme, every item survives correction (Table~\ref{tab:holm_aware}).

\section{Group Comparison by Experience Level}

The four experience cohorts were compared with the Kruskal--Wallis test~\parencite{Kruskal_Wallis_1952}, the rank-based generalization of one-way analysis of variance,
\begin{equation}
  H = \frac{12}{N(N+1)}\sum_{j=1}^{g}\frac{R_j^2}{n_j}-3(N+1),
\end{equation}
with the effect size $\eta^2=(H-g+1)/(N-g)$ interpreted on the small/medium/large scale of \textcite{Cohen_1988}.

\paragraph{Worked example (Awareness, $\eta^2$).} With the omnibus statistic $H=213.096$ across $g=4$ groups and $N=1{,}411$,
\begin{equation}
  \eta^2 = \frac{H-g+1}{N-g} = \frac{213.096-4+1}{1411-4} = \frac{210.096}{1407} = 0.149,
\end{equation}
a medium-to-large effect (Table~\ref{tab:exp}). Post-hoc pairwise contrasts used Dunn's test with Holm adjustment~\parencite{Dunn_1964}, comparing standardized differences of mean ranks with a tie correction. The contrasts localize the effect in the $11$--$15$ year cohort, which scores below both the $1$--$5$ and $6$--$10$ year cohorts; on the Awareness composite worked here, the $16+$ cohort ($n=18$) yields no significant pairwise contrast after correction ($p=.118$ against the $11$--$15$ year cohort), though it does separate significantly from that cohort on the two value composites, as reported in Chapter~4.

\section{Association Analysis}

Associations between categorical variables were tested with the Pearson chi-square test of independence~\parencite{Pearson_1900}, comparing observed counts $O_{ij}$ against the counts expected under independence $E_{ij}=R_iC_j/N$,
\begin{equation}
  \chi^2 = \sum_{i,j}\frac{\left(O_{ij}-E_{ij}\right)^2}{E_{ij}},
  \qquad
  V = \sqrt{\frac{\chi^2}{N\left(\min(r,c)-1\right)}} .
\end{equation}
The second expression is Cram\'er's $V$~\parencite{Cramer_1946}, which rescales $\chi^2$ to a $0$--$1$ effect size.

\paragraph{Worked example (professional field $\times$ prior exposure).} Restricting to the three fields with exposure data, the observed counts ($O$) and the counts expected under independence ($E=R_iC_j/N$) are given below; row totals reflect available-case analysis of the items crossed, so the Engineering row totals $337$ rather than the group size of $339$ because two Engineering respondents did not answer the exposure item Q1.2.
\begin{center}\small
\begin{tabular}{lccc}
\toprule
 & Exposure: Yes & Exposure: No & Row total \\
\midrule
IA/Cyber & $59$ \;($E=113.4$) & $442$ \;($E=387.6$) & $501$ \\
IT       & $161$ \;($E=113.4$) & $340$ \;($E=387.6$) & $501$ \\
Engineering & $83$ \;($E=76.3$) & $254$ \;($E=260.7$) & $337$ \\
\midrule
Column total & $303$ & $1036$ & $N=1339$ \\
\bottomrule
\end{tabular}
\end{center}
Summing $(O-E)^2/E$ over the six cells, for example $(59-113.4)^2/113.4 = 26.1$ and $(161-113.4)^2/113.4 = 20.0$, yields $\chi^2 = 60.334$ on $(3-1)(2-1)=2$ degrees of freedom. The effect size is therefore
\begin{equation}
  V = \sqrt{\frac{\chi^2}{N\left(\min(r,c)-1\right)}}
    = \sqrt{\frac{60.334}{1339\,(2-1)}} = \sqrt{0.0451} = 0.212,
\end{equation}
a moderate association (Table~\ref{tab:assoc}). Experience-by-investment associations follow the same computation.

\section{Bias Diagnostics}

The analysis computed two planned diagnostics, together with a post-hoc straight-lining check. The wave (non-response) analysis ordered the sample by submission time (the timestamp at which each response was submitted) and compared the early- and late-arriving halves with the Mann--Whitney test and Cohen's $d$, implementing the early-versus-late split of the non-response-analysis tradition~\parencite{Armstrong_Overton_1977}. Differences are negligible on every composite (IT~Value $d=-0.02$, early $\bar{x}=4.321$, late $\bar{x}=4.328$; Awareness $d=-0.04$; Security~Value $d=-0.04$; all Mann--Whitney $p>.5$), as Table~\ref{tab:wave} shows, indicating no detectable early-versus-late self-selection in the recruited sample. The priming diagnostic compared value composites between high-familiarity (Q1.1 $\in\{4,5\}$) and low-familiarity (Q1.1 $\in\{1,2\}$) respondents, again with the Mann--Whitney test and Cohen's $d$. Both strata diverge sharply on both composites (IT~Value $\Delta=0.539$, $d=1.30$; Security~Value $\Delta=0.573$, $d=1.25$; Table~\ref{tab:priming}), which weighs against a pure question-framing explanation for the high endorsement observed. Each Cohen's $d$ here is computed by the pooled-standard-deviation formula worked above.

\paragraph{Straight-lining (acquiescence) check.} To gauge the contribution of non-differentiating responding to the high agreement rates, the analysis computed the within-respondent standard deviation across the twenty-two agreement-scale items (Sections~2--5) for each of the $1{,}390$ respondents with the complete battery. Of these, $22$ ($1.6\%$) returned an identical response to all twenty-two items (within-respondent $\mathrm{SD}=0$), and $79$ ($5.7\%$, cumulative) had near-flat response strings ($\mathrm{SD}\le0.25$), i.e., fifty-seven additional respondents; the median within-respondent standard deviation is $0.50$. Only a small minority therefore straight-lined the agreement battery, so acquiescence is a negligible presence in these data and does not manufacture the high endorsement.

\section{Multivariate Model and Robustness}

To separate the contributions of professional field, experience, and familiarity to perceived value, each value composite was regressed on those predictors by ordinary least squares~\parencite{Cohen_Cohen_West_Aiken_2003}, with Information Assurance/Cyber Security and the $1$--$5$ year cohort as reference categories,
\begin{equation}
  \widehat{\text{Value}} = \beta_0 + \beta_1\,\text{Familiarity} + \bm{\beta}_F\,\text{Field} + \bm{\beta}_E\,\text{Experience}.
\end{equation}
The models were fit on the full analytic sample of $N=1{,}411$ and account for a substantial share of variance (IT~Value $R^2=0.64$; Security~Value $R^2=0.67$; Table~\ref{tab:regression}). Three results follow: professional field remains a strong predictor after experience and familiarity are held constant (Information Technology vs.\ Information Assurance $b=0.46$ for IT~Value and $0.50$ for Security~Value, both $p<.001$); the $11$--$15$ year cohort retains a substantial negative coefficient once field and familiarity are controlled ($b=-0.40$, $p<.001$ for IT~Value; $b=-0.53$, $p<.001$ for Security~Value), so the mid-career dip persists as an experience effect net of field composition rather than dissolving into it; and familiarity is a positive predictor of perceived value ($b=0.13$ and $0.14$, both $p<.001$).

\paragraph{Heteroscedasticity-robust refit.} Because the value composites are bounded on the $1$--$5$ scale and left-skewed, the models were refit with HC3 heteroscedasticity-consistent standard errors~\parencite{Long_Ervin_2000}. Every coefficient retained its significance and direction, including the substantial negative $11$--$15$ year coefficient (HC3 $p<.001$ for both value composites), so the conclusions do not depend on the homoscedastic ordinary-least-squares assumption.

\paragraph{Standardized familiarity coefficient.} To express the familiarity effect on a scale-free basis, the outcome and the familiarity predictor were standardized to zero mean and unit variance and the models refit. The resulting standardized coefficient, equivalently $\beta = b\,(s_{\text{Fam}}/s_{\text{Value}})$, gives the change in perceived value, in standard-deviation units, per one-standard-deviation increase in familiarity ($\beta = 0.34$ for IT~Value, $\beta = 0.33$ for Security~Value; both $p<.001$), the values quoted in Chapter~4.

\section{Rank-Order Correlation}

The monotonic association among familiarity, the Awareness composite, and the value composites was summarized with Spearman's rank-order correlation~\parencite{Spearman_1904}, which is the Pearson correlation computed on the ranks of the two variables and therefore assumes neither linearity nor normality,
\begin{equation}
  \rho = 1 - \frac{6\sum_{i=1}^{n} d_i^2}{n\left(n^2-1\right)},
\end{equation}
where $d_i$ is the difference between the ranks of the paired observations for case $i$ (a tie correction is applied when observations share a rank). Familiarity correlates moderately with perceived value ($\rho = 0.632$ with IT~Value, $\rho = 0.616$ with Security~Value; both $p<.001$), whereas the Awareness composite correlates more strongly with both ($\rho = 0.734$ and $\rho = 0.731$). The Awareness composite shares no items with either value composite, so these stronger correlations reflect shared endorsement-style variance rather than substance; they are recorded here for completeness only and are not reported as substantive results in Chapter~4, consistent with the construct-validity assessment later in this appendix.

\section{Supplementary Post-Hoc Robustness and Validity Analyses}\label{sec:appb-posthoc}

The three analyses in this section were not specified in the approved proposal; they were added after data collection to test the robustness and measurement properties of the Chapter~4 findings and are reported as corroboration rather than as confirmatory tests. Each is defined here and carried through to its numerical result.

\subsection{Common-Method Variance: Harman's Single-Factor Test}

Because every measure is a single self-report collected on one occasion, common-method variance could inflate the observed correlations~\parencite{Podsakoff_2003}. Harman's test inspects the unrotated eigenstructure of the $23\times23$ correlation matrix $R$ of the composite items and asks whether one factor explains the majority of the common variance. Writing $\lambda_1\ge\lambda_2\ge\cdots\ge\lambda_{23}$ for the eigenvalues of $R$, whose trace is fixed at $23$, the variance share of the first unrotated factor is
\begin{equation}
  \pi_1 = \frac{\lambda_1}{\sum_{j=1}^{23}\lambda_j} = \frac{\lambda_1}{23}.
\end{equation}
For the analytic battery ($n=1{,}390$ complete cases), two eigenvalues exceed unity and $\lambda_1=10.12$, so $\pi_1 = 10.12/23 = 0.440$. A first factor accounting for $44.0\%$ of the variance, below the $50\%$ dominance heuristic, indicates that no single method factor dominates the data. The test is a lower bound rather than a decisive control, so the conclusion is correspondingly narrow.

\subsection{Construct Validity: Composite Reliability and Average Variance Extracted}

A confirmatory three-factor model, in which each composite is a latent variable indicated by its own items, yields standardized loadings $\lambda_i$ from which composite reliability (CR) and average variance extracted (AVE) follow~\parencite{Fornell_Larcker_1981},
\begin{equation}
  \mathrm{CR} = \frac{\left(\sum_{i} \lambda_i\right)^2}{\left(\sum_{i} \lambda_i\right)^2 + \sum_{i} \left(1-\lambda_i^2\right)},
  \qquad
  \mathrm{AVE} = \frac{1}{k}\sum_{i=1}^{k}\lambda_i^2 .
\end{equation}
Reliability is high for all three composites (CR $=0.86$, $0.87$, and $0.87$), corroborating the Cronbach's $\alpha$ values of the internal-consistency section. Convergent validity, indexed by AVE, falls just short of the $0.50$ benchmark for all three composites ($0.47$, $0.46$, and $0.45$), reflecting item sets that each span a heterogeneous range of capabilities. Discriminant validity is assessed by the Fornell--Larcker criterion, which requires $\sqrt{\mathrm{AVE}}$ for each construct to exceed its correlation with the others. It is not met: the inter-composite correlations ($0.74$ to $0.85$) exceed the corresponding $\sqrt{\mathrm{AVE}}$ ($0.67$ to $0.69$). Model fit itself is good (CFI $=0.966$, RMSEA $=0.040$), so the failure is one of separability, not of fit. The heterotrait--monotrait ratio of correlations~\parencite{Henseler_Ringle_Sarstedt_2015} localizes it: Awareness is discriminant from both value composites (HTMT $=0.83$ against each, below the $0.85$ threshold), whereas IT~Value and Security~Value are not discriminant from each other (HTMT $=0.99$). Both value dimensions are therefore best read as one perceived-value construct measured in two domains, and all three composites are used as reliable summative indices rather than as discriminant latent variables; the group comparisons rest on the reliability of the indices, not on their separability.

One further statistic bears on the internal structure of the Awareness composite, which combines a familiarity self-report with six capability-agreement items. A corrected item-total correlation relates item $y_j$ to the mean of the remaining $k-1$ items. For the Awareness composite it is $0.315$ for Q1.1 against a range of $0.53$ to $0.67$ for the six Section~2 items, and Q1.1 correlates with the mean of those six items at $r=0.32$. Removing Q1.1 raises Cronbach's $\alpha$ from $0.792$ to $0.863$, so the item satisfies both conventional deletion criteria at once, namely a corrected item-total correlation below roughly $0.40$ and an $\alpha$ that improves on removal. All seven items are retained nonetheless, because Chapter~3 specified the composite in advance of data collection, and the AVE of $0.47$ reported above reflects that retained heterogeneity. Chapter~4 records the interpretive consequence: the composite tracks capability endorsement more closely than familiarity, and Q1.1 is reported separately wherever the familiarity finding is at issue.

\subsection{Measurement Invariance Across Professional Field}

To test whether the large Information Technology versus Information Assurance gap could reflect differential item functioning rather than a true difference, each item was regressed on the respondent's own standardized composite score $m$, a field indicator $g$ (with $g=1$ for Information Assurance), and their product, adapting the logistic-regression DIF framework~\parencite{Zumbo_1999} to a linear implementation on the five-point item scores. Uniform DIF is the partial effect of $g$ after matching on $m$, assessed by the improvement in fit of
\begin{equation}
  y_i = \beta_0 + \beta_1 m_i + \beta_2 g_i + \varepsilon_i
\end{equation}
over the matched-only model $y_i=\beta_0+\beta_1 m_i+\varepsilon_i$, with $\beta_2$ giving the Information Assurance versus Information Technology difference in item response, in scale points, at a common trait level. Of the $23$ items, $21$ show no meaningful differential functioning; two (Questions~1.1 and 2.5) meet the flag criterion ($p<.001$ and $|\beta_2|\ge0.20$), the larger effect falling on the familiarity item Q1.1 ($\beta_2=-0.99$ scale points) and a small effect on Q2.5 ($\beta_2=+0.23$). Both flagged items belong to the Awareness composite, so neither value composite contains a differentially functioning item. The large Information Technology versus Information Assurance gaps on the value composites, which approach two standard deviations, therefore cannot be attributed to item-level measurement differences. The same regression's interaction term $m_i g_i$ captures non-uniform DIF, a field difference in how strongly an item tracks the trait. No value-composite item reached the flag threshold on it either, the largest value-composite non-uniform effect falling on Q4.4 at $0.17$ scale points, so the conclusion holds under both uniform and non-uniform screening. That evidence supports reading the field contrast as substantive.

\subsection{Familiarity-Stratum Adjustment of the Field Gap}

Familiarity is bimodal: the two lowest categories hold 78.2\% of respondents and the middle category only 4.2\%, while ``Very familiar'' recovers to 13.6\%. That shape suggests two strata rather than one graded continuum, and it raises the possibility that the professional-field gap reflects unequal draw from those strata rather than field membership itself. Both focal groups draw from these strata, and they do so unequally: 206 of 501 Information Technology respondents fall in the high-familiarity stratum ($\text{Q1.1}\ge4$) against 11 of 501 Information Assurance respondents.

A binary stratum indicator $s_i$ therefore enters the field-only model estimated on the two focal groups ($n=1{,}002$),
\begin{equation}
  y_i = \beta_0 + \beta_1 f_i + \beta_2 s_i + \varepsilon_i,
\end{equation}
where $f_i=1$ for Information Technology and $s_i=1$ for the high-familiarity stratum, with HC3-robust standard errors. Field coefficients fall from $0.689$ to $0.602$ on IT Value and from $0.746$ to $0.655$ on Security Value, retaining $87\%$ of their magnitude on both composites and remaining significant at $p<.001$. Entering the recent-exposure indicator alongside the stratum indicator retains $86\%$ ($b=0.589$ and $b=0.642$ respectively). Stratum composition therefore accounts for roughly one seventh of the field gap.

Observed strata serve here in place of a fitted latent class. A finite-mixture model was considered and set aside: on a five-point ordinal item with a binary covariate, mixture likelihoods admit degenerate components with near-zero variance, and the transparent threshold split reproduces the same partition without that instability. That threshold matches the high-familiarity group already used for the priming diagnostic in Chapter~4, so the two analyses partition the sample identically. This adjustment narrows the recruitment-frame concern without resolving it, since a channel difference that delivered systematically different respondents within a single familiarity stratum would remain invisible to it.

\subsection{Response-Authenticity Screen: Operating Characteristics}\label{sec:appb-forensic}

The response-authenticity screen described in Chapter~4 (Section~\ref{sec:forensic-validation}) was validated against synthetic data with known ground truth, following standard classifier-evaluation practice~\parencite{Fawcett_ROC_2006}, and the operating characteristics are reported so as to separate what the validation confirms from what it genuinely measures. Against contamination defined in the screen's own terms, exact-duplicate and near-duplicate clusters, sensitivity is at or near $100$ percent, but this is a confirmatory check that the rule matches duplicates, not a measure of real-world reach. The informative test plants coordinated fraud from an independent behavioral model, a ring whose members share an answer key but differ in discipline. Detection is then a decay curve: a disciplined ring that shares a demographic profile is caught in full, sensitivity falls as the ring grows sloppier, and it collapses toward zero for a ring whose members each vary their demographic profile, because the screen requires a repeated whole vector. The screen is therefore effective against disciplined, demographically consistent coordinated duplication and, by construction, blind to loosely coordinated, random-responding, or fluent machine-generated submissions, which a complementary content-only detector would be needed to catch.

Specificity was measured against the case most likely to trip the screen: authentic data with genuine ceiling collisions, where strong endorsers land on the same high-agreement profile. On ordinary synthetic data the screen wrongly flags fewer than one authentic row in seven hundred, but on ceiling-heavy data the mechanical rule over-excludes genuine strong endorsers at roughly one to two percent. That rate is close to the borderline-flag rate seen in this study's own data: the $31$ respondents that separate the $N=1{,}411$ primary sample from the $N=1{,}380$ fully screened sample. This is the empirical reason the screen is paired with human adjudication rather than run alone: the automated rule cannot distinguish a genuine ceiling coincidence from a small scripted cluster without the timing and context a reviewer applies. As a further check on real rather than synthetic data, the screen also ran on two large public Likert datasets that bracket this instrument in length. At the twenty-two-item length used here, the false-positive rate on genuinely organic data is negligible, about one row in four thousand, although it rises sharply on much shorter or coarser instruments. Default application is therefore appropriate for instruments of this length and requires recalibration for short ones.

A confirmatory battery of $500$ independently generated datasets bounds the false-positive behavior across a wide distribution of dataset sizes and cluster counts. On $100$ fully organic datasets, comprising $47{,}901$ authentic rows, the screen produced only $7$ false positives, a false-positive rate near $0.01$ percent, and returned a completely clean result on $93$ of the $100$ datasets. The reproducibility artifact carries the complete battery and every validation harness, among them the independent ring model, the ceiling-heavy hard case, and the cross-method complementarity experiment.

\section{Summary}

The analysis applied the estimators and tests above exactly as defined, with non-parametric procedures and effect sizes carrying the interpretive weight given the large sample. The worked examples confirm, step by step, that the headline quantities of Chapter~4 follow directly from the raw responses: the familiarity mean of $1.94$, the Awareness reliability of $\alpha=0.792$, the field effect size of $d=1.64$, the experience effect size of $\eta^2=0.149$, and the field-by-exposure association of $V=0.212$. Every figure in Chapter~4 derives from one of these computations, and each computation is reproducible from the response data through the scripted pipeline described in the Reproducibility and Tooling section.

\section{The Independent Analysis, Reproduced: Code and Output}

The methodology above defines each procedure and works a representative case by hand. This section supplies the executable proof, drawn from the independent statistical analysis and its companion computational notebook. The analysis is a fully scripted Python program (NumPy, pandas, SciPy, and statsmodels, with semopy for the confirmatory measurement model used in the post-hoc validity check) that recomputes every value reported in Chapter~4 directly from the raw response file. Each analysis below presents the exact code followed by its verbatim output, so that no reported quantity rests on assertion: each one is reproducible from the response data by running the code shown. The ordering follows the methodology sections above, and the figures these computations produce appear, with alternative-text descriptions, in Chapter~4.

\subsection{Setup and Data Loading}
Loads the response file and maps the raw columns onto the Chapter~3 question numbering.
\noindent\textbf{Code:}\nopagebreak
{\tiny\begin{verbatim}
import os, json, warnings
import numpy as np
import pandas as pd
from scipy import stats
import matplotlib.pyplot as plt
from matplotlib.patches import Patch
warnings.filterwarnings("ignore")
pd.set_option("display.max_columns", 60)
pd.set_option("display.width", 160)

# The notebook lives in Analysis/; the workbook is one directory up.
ROOT = os.path.dirname(os.getcwd()) if os.path.basename(os.getcwd()) == "Analysis" else os.getcwd()
SRC = os.path.join(ROOT, "Survey Responses Final Export.xlsx")   # formerly "...Final Dump.xlsx"
assert os.path.exists(SRC), f"Cannot find {SRC}"
raw = pd.read_excel(SRC)
cols = list(raw.columns)
print("Rows:", len(raw), "| Columns:", len(cols))
\end{verbatim}}
\noindent\textbf{Output:}\nopagebreak
{\scriptsize\begin{verbatim}
Rows: 1411 | Columns: 35
\end{verbatim}}

\subsection{Data Cleaning and Sample}
Decodes each Likert response to its integer score and confirms the missing-data rule; no retained response crosses the twenty-percent exclusion threshold. The analytic sample follows a documented funnel: $1{,}600$ raw submissions, from which a pre-registered quality screen removed $188$ (completion under three minutes, an open interval exceeding five days, or answering at most Section~1), and a post-hoc response-authenticity screen removed one further submission belonging to a coordinated automated cluster (a single identical response vector submitted $71$ times at roughly one-minute completions), giving the genuine $N=1{,}411$. Complete data on all $23$ Likert items are present for $1{,}390$ respondents ($98.5\%$).
\noindent\textbf{Code:}\nopagebreak
{\tiny\begin{verbatim}
LIKERT = ["Q1.1"]+[f"Q2.{i}" for i in range(1,7)]+[f"Q3.{i}" for i in range(1,7)]+\
         [f"Q4.{i}" for i in range(1,5)]+[f"Q5.{i}" for i in range(1,7)]
BINARY  = ["Q1.2","Q6.1","Q6.2"]
TERNARY = ["Q4.5","Q5.7"]

def to_score(v):
    if pd.isna(v): return np.nan
    s = str(v).strip()
    return int(s[0]) if s and s[0].isdigit() else np.nan

for q in LIKERT:
    df[q] = df[q].map(to_score)
for q in BINARY+TERNARY+["Q6.3","Q6.4"]:
    df[q] = df[q].map(lambda v: np.nan if pd.isna(v) else str(v).strip())

N_TOTAL = len(df)
complete_likert = int(df[LIKERT].notna().all(axis=1).sum())
missing = pd.Series({q:int(df[q].isna().sum()) for q in LIKERT+BINARY+TERNARY+["Q6.3","Q6.4"]})
print(f"N total = {N_TOTAL}")
print(f"Complete across all 23 Likert items = {complete_likert} ({100*complete_likert/N_TOTAL:.1f}%)")
print("Max item-level missing =", int(missing.max()))
missing[missing>0]
\end{verbatim}}
\noindent\textbf{Output:}\nopagebreak
{\scriptsize\begin{verbatim}
N total = 1411
Complete across all 23 Likert items = 1390 (98.5%)
Max item-level missing = 6
Q2.2    6
Q2.3    4
Q2.4    4
Q2.5    3
Q2.6    4
Q3.1    1
Q3.2    3
Q3.3    4
Q3.4    4
Q3.5    3
Q3.6    3
Q4.2    4
Q4.3    3
Q4.4    3
Q5.1    2
Q5.2    3
Q5.3    3
Q5.4    3
Q5.5    3
Q5.6    3
Q1.2    2
Q6.2    6
dtype: int64
\end{verbatim}}

\subsection{Item-Level Descriptive Statistics}
Computes the mean, standard deviation, median, interquartile range, agreement rate, and Wilson interval for every Likert item (methodology: \emph{Descriptive Statistics}).
\noindent\textbf{Code:}\nopagebreak
{\tiny\begin{verbatim}
def ci_mean(x, conf=0.95):
    x = pd.Series(x).dropna(); n = len(x)
    if n < 2: return (np.nan, np.nan)
    se = x.std(ddof=1)/np.sqrt(n)
    h = se*stats.t.ppf((1+conf)/2, n-1)
    return (x.mean()-h, x.mean()+h)

def ci_wilson(k, n, conf=0.95):
    if n == 0: return (np.nan, np.nan)
    z = stats.norm.ppf((1+conf)/2); p = k/n; den = 1+z**2/n
    centre = (p+z**2/(2*n))/den
    half = z*np.sqrt(p*(1-p)/n + z**2/(4*n**2))/den
    return (centre-half, centre+half)

rows = []
for q in LIKERT:
    x = df[q].dropna(); n = len(x); k = int((x>=4).sum())
    lo,hi = ci_mean(x); plo,phi = ci_wilson(k,n)
    rows.append({"Item":q, "n":n, "Mean":round(x.mean(),2), "SD":round(x.std(ddof=1),2),
                 "Mdn":x.median(), "IQR":f"{x.quantile(.25):.0f}-{x.quantile(.75):.0f}",
                 "Agree%":round(100*k/n,1), "CI_low":round(100*plo,1), "CI_high":round(100*phi,1)})
desc = pd.DataFrame(rows).set_index("Item")
desc
\end{verbatim}}
\noindent\textbf{Output:}\nopagebreak
{\scriptsize\begin{verbatim}
         n  Mean    SD  Mdn  IQR  Agree%  CI_low  CI_high
Item
Q1.1  1411  1.94  1.20  2.0  1-2    17.6    15.7     19.6
Q2.1  1411  4.42  0.57  4.0  4-5    95.7    94.6     96.7
Q2.2  1405  4.38  0.56  4.0  4-5    95.8    94.6     96.7
Q2.3  1407  4.39  0.57  4.0  4-5    95.8    94.6     96.7
Q2.4  1407  4.50  0.58  5.0  4-5    95.7    94.5     96.7
Q2.5  1408  4.29  0.72  4.0  4-5    84.6    82.6     86.4
Q2.6  1407  4.41  0.57  4.0  4-5    95.7    94.5     96.7
Q3.1  1410  4.27  0.72  4.0  4-5    84.6    82.6     86.4
Q3.2  1408  4.39  0.57  4.0  4-5    95.7    94.6     96.7
Q3.3  1407  4.51  0.58  5.0  4-5    95.8    94.6     96.7
Q3.4  1407  4.42  0.57  4.0  4-5    95.9    94.7     96.8
Q3.5  1408  4.30  0.72  4.0  4-5    84.8    82.8     86.6
Q3.6  1408  3.89  0.81  4.0  3-5    61.1    58.5     63.6
Q4.1  1411  4.39  0.57  4.0  4-5    95.9    94.7     96.8
Q4.2  1407  4.29  0.71  4.0  4-5    85.5    83.6     87.2
Q4.3  1408  4.30  0.72  4.0  4-5    84.8    82.8     86.6
Q4.4  1408  4.12  0.64  4.0  4-5    84.7    82.7     86.4
Q5.1  1409  4.12  0.64  4.0  4-5    85.0    83.0     86.7
Q5.2  1408  4.13  0.80  4.0  3-5    73.8    71.4     76.0
Q5.3  1408  4.31  0.72  4.0  4-5    84.8    82.8     86.6
Q5.4  1408  4.32  0.75  4.0  4-5    83.0    81.0     84.9
Q5.5  1408  4.22  0.71  4.0  4-5    83.3    81.3     85.2
Q5.6  1408  4.23  0.71  4.0  4-5    83.5    81.5     85.4
\end{verbatim}}

\subsection{Composites and Cronbach's Alpha}
Forms the three composites and computes Cronbach's $\alpha$ for each (methodology: \emph{Composite Scores and Internal Consistency}).
\noindent\textbf{Code:}\nopagebreak
{\scriptsize\begin{verbatim}
composites = {
    "Awareness":  ["Q1.1"]+[f"Q2.{i}" for i in range(1,7)],
    "IT_Value":   ["Q3.1","Q3.3","Q3.4","Q3.5","Q4.1","Q4.2","Q4.3","Q4.4"],
    "Sec_Value":  ["Q3.2","Q3.6","Q5.1","Q5.2","Q5.3","Q5.4","Q5.5","Q5.6"],
}
def cronbach_alpha(frame):
    f = frame.dropna(); k = f.shape[1]
    item_var = f.var(axis=0, ddof=1).sum()
    total_var = f.sum(axis=1).var(ddof=1)
    return (k/(k-1))*(1-item_var/total_var), len(f)

rows = []
for name, items in composites.items():
    s = df[items].mean(axis=1); df[f"COMP_{name}"] = s
    a, na = cronbach_alpha(df[items]); lo,hi = ci_mean(s)
    rows.append({"Composite":name, "k":len(items), "Mean":round(s.mean(),3),
                 "SD":round(s.std(ddof=1),3), "Median":round(s.median(),3),
                 "CI":f"[{lo:.3f}, {hi:.3f}]", "alpha":round(a,3)})
pd.DataFrame(rows).set_index("Composite")
\end{verbatim}}
\noindent\textbf{Output:}\nopagebreak
{\scriptsize\begin{verbatim}
           k   Mean     SD  Median              CI  alpha
Composite                                                
Awareness  7  4.044  0.478   4.000  [4.019, 4.069]  0.792
IT_Value   8  4.325  0.463   4.375  [4.301, 4.349]  0.858
Sec_Value  8  4.202  0.509   4.250  [4.176, 4.229]  0.858
\end{verbatim}}

\subsection{Shapiro--Wilk Normality Tests}
Applies the Shapiro--Wilk normality test to each composite (methodology: \emph{Tests of Normality}).
\noindent\textbf{Code:}\nopagebreak
{\scriptsize\begin{verbatim}
for name in composites:
    s = df[f"COMP_{name}"].dropna()
    W,p = stats.shapiro(s.sample(min(5000,len(s)), random_state=1))
    print(f"{name:10s}  W={W:.4f}  p={p:.2e}  -> {'non-normal' if p<0.05 else 'normal'}")
\end{verbatim}}
\noindent\textbf{Output:}\nopagebreak
{\scriptsize\begin{verbatim}
Awareness   W=0.9780  p=7.49e-14  -> non-normal
IT_Value    W=0.9571  p=5.93e-20  -> non-normal
Sec_Value   W=0.9647  p=4.80e-18  -> non-normal
\end{verbatim}}

\subsection{Effect-Size Estimators}
Defines the effect-size estimators: Cohen's $d$, the rank-biserial correlation, and the large-sample confidence interval for $d$ (methodology: \emph{Effect Sizes}).
\noindent\textbf{Code:}\nopagebreak
{\scriptsize\begin{verbatim}
def cohen_d(a,b):
    a,b = np.asarray(a), np.asarray(b); na,nb = len(a),len(b)
    sp = np.sqrt(((na-1)*a.var(ddof=1)+(nb-1)*b.var(ddof=1))/(na+nb-2))
    return (a.mean()-b.mean())/sp if sp>0 else np.nan
def rank_biserial(U,n1,n2):
    return 1 - 2*U/(n1*n2)
def d_ci(d, n1, n2):
    """Large-sample 95% CI for Cohen's d (Hedges & Olkin, 1985)."""
    se = np.sqrt((n1+n2)/(n1*n2) + d**2/(2*(n1+n2)))
    return d-1.96*se, d+1.96*se
print("Effect-size helpers defined.")
\end{verbatim}}
\noindent\textbf{Output:}\nopagebreak
{\scriptsize\begin{verbatim}
Effect-size helpers defined.
\end{verbatim}}

\subsection{Field Comparison: Mann--Whitney, Cohen's d, Rank-Biserial}
Compares Information Technology and Information Assurance on each composite with the Mann--Whitney $U$ test, Cohen's $d$ with its interval, and the rank-biserial correlation (methodology: \emph{Group Comparison by Professional Field}).
\noindent\textbf{Code:}\nopagebreak
{\tiny\begin{verbatim}
it_mask = df["Q6.3"]=="Information Technology"
ia_mask = df["Q6.3"]=="Information Assurance / Cyber Security"
rows = []
for name in composites:
    a = df.loc[it_mask, f"COMP_{name}"].dropna()
    b = df.loc[ia_mask, f"COMP_{name}"].dropna()
    U,p = stats.mannwhitneyu(a,b,alternative="two-sided")
    t,tp = stats.ttest_ind(a,b,equal_var=False)
    d = round(cohen_d(a,b),3); lo,hi = d_ci(d,len(a),len(b))  # round d as in the published table pipeline
    rows.append({"Composite":name,"IT n":len(a),"IA n":len(b),
                 "IT M":round(a.mean(),3),"IA M":round(b.mean(),3),
                 "U":U,"p":f"{p:.2e}","Cohen d":round(d,3),
                 "d 95% CI":f"[{lo:.2f}, {hi:.2f}]",
                 "r_rb":round(rank_biserial(U,len(a),len(b)),3)})
pd.DataFrame(rows).set_index("Composite")
\end{verbatim}}
\noindent\textbf{Output:}\nopagebreak
{\scriptsize\begin{verbatim}
           IT n  IA n   IT M   IA M         U          p  Cohen d      d 95% CI   r_rb
Composite
Awareness   501   501  4.379  3.732  218765.5   8.46e-93    1.641  [1.50, 1.78] -0.743
IT_Value    501   501  4.669  3.980  227984.0  1.29e-111    1.940  [1.79, 2.09] -0.817
Sec_Value   501   501  4.566  3.820  227813.5  4.60e-111    1.944  [1.79, 2.09] -0.815
\end{verbatim}}

\subsection{Item-Level Comparison with Holm--Bonferroni}
Repeats the comparison item by item and applies the Holm--Bonferroni step-down correction within each item family.
\noindent\textbf{Code:}\nopagebreak
{\tiny\begin{verbatim}
families = {
    "Awareness (S1-S2)": ["Q1.1"]+[f"Q2.{i}" for i in range(1,7)],
    "Applicability (S3)": [f"Q3.{i}" for i in range(1,7)],
    "IT Value (S4)": [f"Q4.{i}" for i in range(1,5)],
    "IA Value (S5)": [f"Q5.{i}" for i in range(1,7)],
}
def holm(pairs):
    m = len(pairs); order = sorted(pairs, key=lambda kv: kv[1])
    out = {}; prev = 0.0
    for i,(k,p) in enumerate(order):
        adj = min(1.0,(m-i)*p); adj = max(adj,prev); prev = adj
        out[k] = {"p":p,"p_adj":adj,"reject":adj<0.05}
    return out

allrows = []
for fam, items in families.items():
    pairs = []; cache = {}
    for q in items:
        a = df.loc[it_mask,q].dropna(); b = df.loc[ia_mask,q].dropna()
        U,p = stats.mannwhitneyu(a,b,alternative="two-sided")
        cache[q] = {"IT M":round(a.mean(),2),"IA M":round(b.mean(),2),
                    "r_rb":round(rank_biserial(U,len(a),len(b)),3)}
        pairs.append((q,p))
    h = holm(pairs)
    for q in items:
        allrows.append({"Family":fam,"Item":q,**cache[q],
                        "p_Holm":f"{h[q]['p_adj']:.1e}","Sig":"yes" if h[q]['reject'] else "no"})
pd.DataFrame(allrows).set_index(["Family","Item"])
\end{verbatim}}
\noindent\textbf{Output:}\nopagebreak
{\scriptsize\begin{verbatim}
                         IT M  IA M   r_rb    p_Holm  Sig
Family             Item
Awareness (S1-S2)  Q1.1  2.88  1.20 -0.775  6.4e-115  yes
                   Q2.1  4.65  4.17 -0.415   5.8e-37  yes
                   Q2.2  4.60  4.15 -0.395   4.3e-34  yes
                   Q2.3  4.61  4.16 -0.398   2.1e-34  yes
                   Q2.4  4.73  4.25 -0.416   1.3e-37  yes
                   Q2.5  4.53  4.02 -0.365   2.1e-27  yes
                   Q2.6  4.65  4.17 -0.428   3.6e-39  yes
Applicability (S3) Q3.1  4.62  3.93 -0.519   2.9e-53  yes
                   Q3.2  4.74  4.05 -0.613   1.0e-79  yes
                   Q3.3  4.85  4.13 -0.634   1.0e-85  yes
                   Q3.4  4.80  4.05 -0.660   2.4e-91  yes
                   Q3.5  4.64  3.97 -0.494   9.1e-49  yes
                   Q3.6  4.21  3.54 -0.448   5.5e-39  yes
IT Value (S4)      Q4.1  4.73  4.03 -0.621   1.1e-81  yes
                   Q4.2  4.61  3.93 -0.503   2.5e-50  yes
                   Q4.3  4.69  3.94 -0.553   7.0e-61  yes
                   Q4.4  4.42  3.85 -0.445   1.0e-41  yes
IA Value (S5)      Q5.1  4.47  3.76 -0.564   5.1e-66  yes
                   Q5.2  4.43  3.78 -0.444   2.3e-38  yes
                   Q5.3  4.70  3.92 -0.573   5.4e-65  yes
                   Q5.4  4.69  3.95 -0.525   1.6e-55  yes
                   Q5.5  4.64  3.77 -0.637   2.6e-78  yes
                   Q5.6  4.64  3.78 -0.639   9.8e-79  yes
\end{verbatim}}

\subsection{Experience: Kruskal--Wallis, Eta-Squared, and Dunn}
Compares the four experience cohorts with the Kruskal--Wallis test, computes $\eta^2$, and runs Dunn's post-hoc test with Holm adjustment (methodology: \emph{Group Comparison by Experience Level}).
\noindent\textbf{Code:}\nopagebreak
{\tiny\begin{verbatim}
EXP = ["1-5 Years","6-10 Years","11-15 Years","16+ Years"]
def dunn(groups, names):
    allv = np.concatenate(groups); ranks = stats.rankdata(allv); N=len(allv)
    idx=0; rg=[]
    for g in groups: rg.append(ranks[idx:idx+len(g)]); idx+=len(g)
    Rbar=[r.mean() for r in rg]; ns=[len(g) for g in groups]
    _,counts = np.unique(allv, return_counts=True)
    ties=(counts**3-counts).sum(); sigma2=(N*(N+1)/12)-ties/(12*(N-1))
    pairs=[]
    for i in range(len(groups)):
        for j in range(i+1,len(groups)):
            se=np.sqrt(sigma2*(1/ns[i]+1/ns[j])); z=(Rbar[i]-Rbar[j])/se
            p=2*(1-stats.norm.cdf(abs(z))); pairs.append((f"{names[i]} vs {names[j]}",p,round(z,2)))
    h=holm([(k,p) for k,p,_ in pairs])
    return pd.DataFrame([{"pair":k,"z":z,"p_adj":f"{h[k]['p_adj']:.2e}","sig":h[k]['reject']} for k,p,z in pairs])

for name in composites:
    groups=[df.loc[df["Q6.4"]==l,f"COMP_{name}"].dropna().values for l in EXP]
    H,p = stats.kruskal(*groups); Ntot=sum(len(g) for g in groups)
    eta2=(H-len(groups)+1)/(Ntot-len(groups))
    print(f"=== {name} ===  H={H:.1f}  df={len(groups)-1}  p={p:.2e}  eta^2={eta2:.3f}")
    print({l:round(g.mean(),3) for l,g in zip(EXP,groups)})
    display(dunn(groups,EXP))
\end{verbatim}}
\noindent\textbf{Output:}\nopagebreak
{\tiny\begin{verbatim}
=== Awareness ===  H=213.1  df=3  p=6.24e-46  eta^2=0.149
{'1-5 Years': np.float64(4.063), '6-10 Years': np.float64(4.236), '11-15 Years': np.float64(3.735), '16+ Years': np.float64(4.008)}
                        pair      z     p_adj    sig
0    1-5 Years vs 6-10 Years  -6.09  4.58e-09   True
1   1-5 Years vs 11-15 Years   9.51  0.00e+00   True
2     1-5 Years vs 16+ Years   0.82  4.11e-01  False
3  6-10 Years vs 11-15 Years  14.56  0.00e+00   True
4    6-10 Years vs 16+ Years   2.39  5.11e-02  False
5   11-15 Years vs 16+ Years  -1.89  1.18e-01  False=== IT_Value ===  H=229.4  df=3  p=1.88e-49  eta^2=0.161
{'1-5 Years': np.float64(4.402), '6-10 Years': np.float64(4.462), '11-15 Years': np.float64(3.982), '16+ Years': np.float64(4.59)}
                        pair      z     p_adj    sig
0    1-5 Years vs 6-10 Years  -2.20  8.42e-02  False
1   1-5 Years vs 11-15 Years  12.52  0.00e+00   True
2     1-5 Years vs 16+ Years  -1.79  1.46e-01  False
3  6-10 Years vs 11-15 Years  14.09  0.00e+00   True
4    6-10 Years vs 16+ Years  -1.22  2.21e-01  False
5   11-15 Years vs 16+ Years  -5.33  3.88e-07   True=== Sec_Value ===  H=291.2  df=3  p=7.80e-63  eta^2=0.205
{'1-5 Years': np.float64(4.323), '6-10 Years': np.float64(4.352), '11-15 Years': np.float64(3.769), '16+ Years': np.float64(4.431)}
                        pair      z     p_adj    sig
0    1-5 Years vs 6-10 Years  -0.80  1.00e+00  False
1   1-5 Years vs 11-15 Years  15.07  0.00e+00   True
2     1-5 Years vs 16+ Years  -0.84  1.00e+00  False
3  6-10 Years vs 11-15 Years  15.35  0.00e+00   True
4    6-10 Years vs 16+ Years  -0.63  1.00e+00  False
5   11-15 Years vs 16+ Years  -5.12  1.24e-06   True
\end{verbatim}}

\subsection{Associations: Chi-Square and Cram\'er's V}
Tests the categorical associations with the Pearson chi-square test and Cram\'er's $V$ (methodology: \emph{Association Analysis}).
\noindent\textbf{Code:}\nopagebreak
{\scriptsize\begin{verbatim}
def chi_assoc(rowvar, colvar, roworder=None, colorder=None):
    sub = df[[rowvar,colvar]].dropna()
    ct = pd.crosstab(sub[rowvar], sub[colvar])
    if roworder: ct = ct.reindex(roworder).dropna(how="all")
    if colorder: ct = ct.reindex(columns=colorder).dropna(axis=1,how="all")
    chi2,p,dof,_ = stats.chi2_contingency(ct)
    n = ct.values.sum(); k = min(ct.shape)-1
    V = np.sqrt(chi2/(n*k)) if k>0 else np.nan
    return ct, chi2, dof, p, V

field_order=["Information Assurance / Cyber Security","Information Technology","Engineering"]
for lab,(rv,cv,ro,co) in {
    "Field x exposure":("Q6.3","Q1.2",field_order,["Yes","No"]),
    "Experience x IT investment":("Q6.4","Q4.5",EXP,["Yes","No","Unsure"]),
    "Experience x Sec investment":("Q6.4","Q5.7",EXP,["Yes","No","Unsure"]),
}.items():
    ct,chi2,dof,p,V = chi_assoc(rv,cv,ro,co)
    print(f"{lab}: chi2({dof})={chi2:.1f}, p={p:.2e}, Cramer V={V:.3f}, n={ct.values.sum()}")
\end{verbatim}}
\noindent\textbf{Output:}\nopagebreak
{\scriptsize\begin{verbatim}
Field x exposure: chi2(2)=60.3, p=7.92e-14, Cramer V=0.212, n=1339
Experience x IT investment: chi2(6)=6.0, p=4.21e-01, Cramer V=0.046, n=1411
Experience x Sec investment: chi2(6)=2.2, p=9.02e-01, Cramer V=0.028, n=1411
\end{verbatim}}

\subsection{Bias Diagnostics}
Runs the three bias diagnostics: the wave (non-response) analysis, the priming diagnostic, and the straight-lining check (methodology: \emph{Bias Diagnostics}).
\noindent\textbf{Code:}\nopagebreak
{\tiny\begin{verbatim}
# Wave -- "Completion time" is the submission timestamp, so sorting it orders by arrival
ds = df.sort_values("Completion time"); half=len(ds)//2
early, late = ds.iloc[:half], ds.iloc[half:]
print("WAVE (early vs late):")
for name in composites:
    a=early[f"COMP_{name}"].dropna(); b=late[f"COMP_{name}"].dropna()
    U,p=stats.mannwhitneyu(a,b,alternative="two-sided")
    print(f"  {name:10s} early={a.mean():.3f} late={b.mean():.3f} d={cohen_d(a,b):+.3f} p={p:.2e}")
# Priming
high=df[df["Q1.1"].isin([4,5])]; low=df[df["Q1.1"].isin([1,2])]
print("\nPRIMING (high vs low familiarity):")
for name in ["IT_Value","Sec_Value"]:
    a=high[f"COMP_{name}"].dropna(); b=low[f"COMP_{name}"].dropna()
    U,p=stats.mannwhitneyu(a,b,alternative="two-sided")
    print(f"  {name:10s} high={a.mean():.3f} low={b.mean():.3f} diff={a.mean()-b.mean():+.3f} d={cohen_d(a,b):.3f} p={p:.2e}")

# Straight-lining / acquiescence check across the 22 agreement-scale items
AGREE22 = ([f"Q2.{i}" for i in range(1,7)]+[f"Q3.{i}" for i in range(1,7)]
           +[f"Q4.{i}" for i in range(1,5)]+[f"Q5.{i}" for i in range(1,7)])
batt = df[AGREE22]; complete = batt.notna().all(axis=1)
rowsd = batt[complete].std(axis=1, ddof=0)
nC = int(complete.sum()); sd0 = int((rowsd==0).sum()); flat = int((rowsd<=0.25).sum())
print(f"\nSTRAIGHT-LINING (22 agreement items, n={nC} with complete battery):")
print(f"  identical response to all 22 (within-respondent SD=0): {sd0} ({100*sd0/nC:.1f}%)")
print(f"  near-flat (SD <= 0.25):                                {flat} ({100*flat/nC:.1f}%)")
print(f"  median within-respondent SD:                           {rowsd.median():.3f}")
\end{verbatim}}
\noindent\textbf{Output:}\nopagebreak
{\scriptsize\begin{verbatim}
WAVE (early vs late):
  Awareness  early=4.033 late=4.054 d=-0.044 p=5.23e-01
  IT_Value   early=4.321 late=4.328 d=-0.016 p=8.02e-01
  Sec_Value  early=4.192 late=4.212 d=-0.040 p=5.27e-01

PRIMING (high vs low familiarity):
  IT_Value   high=4.751 low=4.212 diff=+0.539 d=1.302 p=2.71e-67
  Sec_Value  high=4.654 low=4.081 diff=+0.573 d=1.253 p=2.78e-61

STRAIGHT-LINING (22 agreement items, n=1390 with complete battery):
  identical response to all 22 (within-respondent SD=0): 22 (1.6%)
  near-flat (SD <= 0.25):                                79 (5.7%)
  median within-respondent SD:                           0.498
\end{verbatim}}

\subsection{Multivariate Regression with HC3 Robustness}
Fits the ordinary-least-squares regression of each value composite on familiarity, field, and experience, refits with HC3-robust standard errors, and reports the coefficient table and standardized betas (methodology: \emph{Multivariate Model and Robustness}).
\noindent\textbf{Code:}\nopagebreak
{\scriptsize\begin{verbatim}
import statsmodels.formula.api as smf

reg = df.copy()
reg["field4"] = pd.Categorical(
    reg["Q6.3"].map({"Information Technology": "IT",
                     "Information Assurance / Cyber Security": "IA",
                     "Engineering": "Eng", "Other": "Other"}),
    categories=["IA", "IT", "Eng", "Other"])          # IA/Cyber = reference
reg["expc"] = pd.Categorical(
    reg["Q6.4"].map({"1-5 Years": "1-5", "6-10 Years": "6-10",
                     "11-15 Years": "11-15", "16+ Years": "16+"}),
    categories=["1-5", "6-10", "11-15", "16+"])        # 1-5 years = reference
reg["Fam"] = reg["Q1.1"]
reg = reg.dropna(subset=["COMP_IT_Value", "COMP_Sec_Value", "field4", "expc", "Fam"])

LABELS = {
    "Intercept": "Intercept",
    "Fam": "Familiarity (Q1.1)",
    "C(field4, Treatment('IA'))[T.IT]":    "Field: Information Technology",
    "C(field4, Treatment('IA'))[T.Eng]":   "Field: Engineering",
    "C(field4, Treatment('IA'))[T.Other]": "Field: Other",
    "C(expc, Treatment('1-5'))[T.6-10]":   "Experience: 6-10 years",
    "C(expc, Treatment('1-5'))[T.11-15]":  "Experience: 11-15 years",
    "C(expc, Treatment('1-5'))[T.16+]":    "Experience: 16+ years",
}
stars = lambda p: "***" if p < .001 else "**" if p < .01 else "*" if p < .05 else ""

cols_out, fits = {}, {}
for col, title in [("COMP_IT_Value", "IT Value"), ("COMP_Sec_Value", "Security Value")]:
    m = smf.ols(f"{col} ~ Fam + C(field4, Treatment('IA')) + C(expc, Treatment('1-5'))",
                data=reg).fit()
    fits[title] = m
    cols_out[title] = {LABELS.get(k, k): f"{v:.3f}{stars(m.pvalues[k])} ({m.bse[k]:.3f})"
                       for k, v in m.params.items()}

order = ["Familiarity (Q1.1)", "Field: Information Technology", "Field: Engineering",
         "Field: Other", "Experience: 6-10 years", "Experience: 11-15 years",
         "Experience: 16+ years", "Intercept"]
tab = pd.DataFrame(cols_out).reindex(order)
tab.loc["R-squared"]      = [f"{fits[t].rsquared:.3f}"     for t in tab.columns]
tab.loc["Adj. R-squared"] = [f"{fits[t].rsquared_adj:.3f}" for t in tab.columns]
tab.loc["N"]              = [f"{int(fits[t].nobs):,}"      for t in tab.columns]
print("Reference categories: Information Assurance / Cyber Security; 1-5 years of experience.")
print("Coefficients unstandardized (SE in parentheses); * p<.05, ** p<.01, *** p<.001.")
print("The 11-15 year coefficient is large, negative, and significant in both models: the")
print("mid-career dip persists net of field and familiarity, a genuine experience effect.")
display(tab)

# HC3 heteroscedasticity-robust refit (bounded, left-skewed composites)
hc3_rows = []
for col, title in [("COMP_IT_Value", "IT Value"), ("COMP_Sec_Value", "Security Value")]:
    m = fits[title]; mh = m.get_robustcov_results(cov_type="HC3")
    names = list(m.params.index)
    for key, lab in [("Fam", "Familiarity (Q1.1)"),
                     ("C(field4, Treatment('IA'))[T.IT]", "Field: Information Technology"),
                     ("C(expc, Treatment('1-5'))[T.11-15]", "Experience: 11-15 years")]:
        j = names.index(key)
        hc3_rows.append({"Model": title, "Predictor": lab,
                         "b": round(float(m.params[key]), 3),
                         "SE (classical)": round(float(m.bse[key]), 3),
                         "p (classical)": round(float(m.pvalues[key]), 3),
                         "SE (HC3)": round(float(mh.bse[j]), 3),
                         "p (HC3)": round(float(mh.pvalues[j]), 3)})
print("\nHC3 robustness check (bounded, left-skewed composites):")
print("Significance and direction are unchanged; the 11-15 coefficient stays significant.")
pd.DataFrame(hc3_rows).set_index(["Model", "Predictor"])
\end{verbatim}}
\noindent\textbf{Output:}\nopagebreak
{\tiny\begin{verbatim}
Reference categories: Information Assurance / Cyber Security; 1-5 years of experience.
Coefficients unstandardized (SE in parentheses); * p<.05, ** p<.01, *** p<.001.
The 11-15 year coefficient is large, negative, and significant in both models: the
mid-career dip persists net of field and familiarity, a genuine experience effect.
                                        IT Value     Security Value
Familiarity (Q1.1)              0.132*** (0.008)   0.142*** (0.008)
Field: Information Technology   0.458*** (0.022)   0.496*** (0.023)
Field: Engineering              0.301*** (0.020)   0.350*** (0.021)
Field: Other                    0.232*** (0.036)   0.305*** (0.038)
Experience: 6-10 years          0.082*** (0.017)    0.055** (0.018)
Experience: 11-15 years        -0.402*** (0.019)  -0.534*** (0.020)
Experience: 16+ years             0.168* (0.067)      0.092 (0.070)
Intercept                       3.888*** (0.018)   3.759*** (0.019)
R-squared                                  0.635              0.672
Adj. R-squared                             0.633              0.670
N                                          1,411              1,411
HC3 robustness check (bounded, left-skewed composites):
Significance and direction are unchanged; the 11-15 coefficient stays significant.
                                                  b  SE (classical)  p (classical)  SE (HC3)  p (HC3)
Model          Predictor
IT Value       Familiarity (Q1.1)             0.132           0.008            0.0     0.007      0.0
               Field: Information Technology  0.458           0.022            0.0     0.021      0.0
               Experience: 11-15 years       -0.402           0.019            0.0     0.020      0.0
Security Value Familiarity (Q1.1)             0.142           0.008            0.0     0.008      0.0
               Field: Information Technology  0.496           0.023            0.0     0.023      0.0
               Experience: 11-15 years       -0.534           0.020            0.0     0.020      0.0
\end{verbatim}}

\subsection{Spearman Rank-Order Correlations}
Computes the Spearman rank-order correlations among familiarity, the Awareness composite, and the value composites (methodology: \emph{Rank-Order Correlation}).
\noindent\textbf{Code:}\nopagebreak
{\scriptsize\begin{verbatim}
from scipy.stats import spearmanr

pairs = [("Q1.1", "COMP_IT_Value", "Familiarity (Q1.1)", "IT Value"),
         ("Q1.1", "COMP_Sec_Value", "Familiarity (Q1.1)", "Security Value"),
         ("COMP_Awareness", "COMP_IT_Value", "Awareness composite", "IT Value"),
         ("COMP_Awareness", "COMP_Sec_Value", "Awareness composite", "Security Value"),
         ("COMP_IT_Value", "COMP_Sec_Value", "IT Value", "Security Value")]
rows = []
for a, b, la, lb in pairs:
    s = df[[a, b]].dropna()
    rho, p = spearmanr(s[a], s[b])
    rows.append({"X": la, "Y": lb, "Spearman rho": round(rho, 3),
                 "p": f"{p:.2e}", "n": len(s)})
pd.DataFrame(rows)
\end{verbatim}}
\noindent\textbf{Output:}\nopagebreak
{\scriptsize\begin{verbatim}
                     X               Y  Spearman rho          p     n
0   Familiarity (Q1.1)        IT Value         0.632  4.53e-158  1411
1   Familiarity (Q1.1)  Security Value         0.616  2.27e-148  1411
2  Awareness composite        IT Value         0.734  3.48e-239  1411
3  Awareness composite  Security Value         0.731  4.90e-236  1411
4             IT Value  Security Value         0.842   0.00e+00  1411
\end{verbatim}}

\subsection{Post-Hoc Common-Method Variance (Harman's Single-Factor Test)}
A robustness check added after data collection and not specified in the proposal. It confirms that no single factor dominates the item covariances.
\noindent\textbf{Code:}\nopagebreak
{\tiny\begin{verbatim}
# Harman's single-factor test on the 23 composite items
Z = (df[LIKERT] - df[LIKERT].mean()) / df[LIKERT].std(ddof=1)
R = np.corrcoef(Z.dropna().values, rowvar=False)
eig = np.sort(np.linalg.eigvalsh(R))[::-1]
print("Items:", len(LIKERT))
print("Factors with eigenvalue > 1:", int((eig > 1).sum()))
print(f"First unrotated factor variance: {100*eig[0]/eig.sum():.1f}%")
\end{verbatim}}
\noindent\textbf{Output:}\nopagebreak
{\scriptsize\begin{verbatim}
Items: 23
Factors with eigenvalue > 1: 2
First unrotated factor variance: 44.0%
\end{verbatim}}

\subsection{Post-Hoc Construct Validity (CR, AVE, Fornell--Larcker)}
A confirmatory three-factor model (semopy) yielding composite reliability, average variance extracted, the discriminant-validity check, and fit indices.
\noindent\textbf{Code:}\nopagebreak
{\tiny\begin{verbatim}
import semopy
san = lambda s: s.replace(".", "_")
X = df[LIKERT].dropna()
desc = "\n".join(f"{n} =~ " + " + ".join(san(i) for i in items)
                 for n, items in composites.items())
cfa = semopy.Model(desc); cfa.fit(X.rename(columns=san))
ins = cfa.inspect(std_est=True)
load = {r["lval"]: abs(float(r["Est. Std"])) for _, r in ins[ins["op"] == "~"].iterrows()}
for n, items in composites.items():
    lam = np.array([load[san(i)] for i in items])
    ave = np.mean(lam**2); cr = lam.sum()**2 / (lam.sum()**2 + np.sum(1 - lam**2))
    print(f"{n:10s} CR={cr:.3f} AVE={ave:.3f} sqrtAVE={ave**0.5:.3f}")
comp = pd.DataFrame({n: df[items].mean(axis=1) for n, items in composites.items()})
print(comp.corr().round(3))
stt = semopy.calc_stats(cfa).T
print(f"CFI={float(stt.loc['CFI'].iloc[0]):.3f} RMSEA={float(stt.loc['RMSEA'].iloc[0]):.3f}")
def htmt(a, b):
    C = df[a + b].corr().abs()
    het = np.mean([C.loc[i, j] for i in a for j in b])
    mono = lambda g: np.mean([C.loc[g[i], g[j]] for i in range(len(g)) for j in range(i+1, len(g))])
    return het / np.sqrt(mono(a) * mono(b))
nm = list(composites); it = list(composites.values())
for i in range(3):
    for j in range(i+1, 3):
        print(f"HTMT {nm[i]}-{nm[j]}: {htmt(it[i], it[j]):.3f}")
\end{verbatim}}
\noindent\textbf{Output:}\nopagebreak
{\scriptsize\begin{verbatim}
Awareness  CR=0.858 AVE=0.471 sqrtAVE=0.686
IT_Value   CR=0.868 AVE=0.455 sqrtAVE=0.675
Sec_Value  CR=0.867 AVE=0.453 sqrtAVE=0.673
           Awareness  IT_Value  Sec_Value
Awareness      1.000     0.740      0.737
IT_Value       0.740     1.000      0.845
Sec_Value      0.737     0.845      1.000
CFI=0.966 RMSEA=0.040
HTMT Awareness-IT_Value: 0.830
HTMT Awareness-Sec_Value: 0.830
HTMT IT_Value-Sec_Value: 0.987
\end{verbatim}}

\subsection{Post-Hoc Measurement Invariance: Linear DIF Screen (Zumbo's framework, linear implementation), IT vs IA}
An item-level differential-item-functioning screen across the two focal fields, matching each item on its composite score.
\noindent\textbf{Code:}\nopagebreak
{\tiny\begin{verbatim}
import statsmodels.api as sm
from scipy.stats import f as _f
foc = df[df["Q6.3"].isin(["Information Technology",
                          "Information Assurance / Cyber Security"])].copy()
foc["grp"] = (foc["Q6.3"] == "Information Assurance / Cyber Security").astype(int)
comp_of = {i: n for n, items in composites.items() for i in items}
compcol = {n: df[items].mean(axis=1) for n, items in composites.items()}
flagged = []
for it in LIKERT:
    sub = pd.DataFrame({"y": foc[it], "m": compcol[comp_of[it]].reindex(foc.index),
                        "g": foc["grp"]}).dropna()
    mm = (sub["m"] - sub["m"].mean()) / sub["m"].std()
    y = sub["y"].values
    r  = sm.OLS(y, sm.add_constant(np.column_stack([mm, sub["g"].values]))).fit()
    m0 = sm.OLS(y, sm.add_constant(mm)).fit()
    F = (m0.ssr - r.ssr) / (r.ssr / (len(y) - 3))
    p = 1 - _f.cdf(F, 1, len(y) - 3)
    if p < 0.001 and abs(r.params[2]) >= 0.20:
        flagged.append((it, round(float(r.params[2]), 3)))
print("Items tested:", len(LIKERT), "| flagged:", len(flagged))
print(flagged)
\end{verbatim}}
\noindent\textbf{Output:}\nopagebreak
{\scriptsize\begin{verbatim}
Items tested: 23 | flagged: 2
[('Q1.1', -0.988), ('Q2.5', 0.23)]
\end{verbatim}}

\subsection{Post-Hoc Awareness-Composite Item Diagnostics and Familiarity-Stratum Adjustment}
Two checks added in revision 1.9. The first reports the corrected item-total correlations and alpha-if-deleted for the Awareness composite, documenting why the familiarity item is a deletion candidate on the conventional criteria and what the pre-specified retention decision costs. The second holds the composition of the bimodal familiarity distribution constant and re-estimates the professional-field gap.
\noindent\textbf{Code:}\nopagebreak
{\tiny\begin{verbatim}
AWARE_Q2 = [f"Q2.{i}" for i in range(1, 7)]
AWARE = ["Q1.1"] + AWARE_Q2

def cronbach(frame):
    k = frame.shape[1]
    return k / (k - 1) * (1 - frame.var(ddof=1).sum() / frame.sum(axis=1).var(ddof=1))

X = df[AWARE].dropna()
print(f"alpha, all seven items : {cronbach(X):.4f}")
print(f"alpha, Q1.1 deleted    : {cronbach(X[AWARE_Q2]):.4f}")
for q in AWARE:
    r = X[q].corr(X[[j for j in AWARE if j != q]].mean(axis=1))
    print(f"  {q:6s} {r:+.3f}")

# familiarity-stratum adjustment of the field gap
IA = "Information Assurance / Cyber Security"
work = df.copy()
for _n, _items in composites.items():
    work[_n] = df[_items].mean(axis=1)
foc = work[work["Q6.3"].isin(["Information Technology", IA])].dropna(
    subset=["Q1.1", "IT_Value", "Sec_Value"]).copy()
foc["IT_field"] = (foc["Q6.3"] == "Information Technology").astype(int)
foc["expo"]  = (foc["Q1.2"] == "Yes").astype(float)
foc["famhi"] = (foc["Q1.1"] >= 4).astype(int)
specs = {"field only": ["IT_field"], "+ exposure (Q1.2)": ["IT_field", "expo"],
         "+ familiarity stratum": ["IT_field", "famhi"],
         "+ both": ["IT_field", "expo", "famhi"]}
for dv in ["IT_Value", "Sec_Value"]:
    base = None
    for name, cols in specs.items():
        m = sm.OLS(foc[dv], sm.add_constant(foc[cols])).fit(cov_type="HC3")
        b = m.params["IT_field"]; base = b if base is None else base
        print(f"{dv} {name:22s} b={b:.3f} retained={100*b/base:5.1f}%")
\end{verbatim}}
\noindent\textbf{Output:}\nopagebreak
{\scriptsize\begin{verbatim}
alpha, all seven items : 0.7923
alpha, Q1.1 deleted    : 0.8627
  Q1.1   +0.315
  Q2.1   +0.655
  Q2.2   +0.613
  Q2.3   +0.618
  Q2.4   +0.671
  Q2.5   +0.532
  Q2.6   +0.671
IT_Value  field only             b=0.689 retained=100.0%
IT_Value  + exposure (Q1.2)      b=0.667 retained= 96.8%
IT_Value  + familiarity stratum  b=0.602 retained= 87.3%
IT_Value  + both                 b=0.589 retained= 85.5%
Sec_Value field only             b=0.746 retained=100.0%
Sec_Value + exposure (Q1.2)      b=0.724 retained= 97.0%
Sec_Value + familiarity stratum  b=0.655 retained= 87.7%
Sec_Value + both                 b=0.642 retained= 86.0%
\end{verbatim}}

\subsection{Synthesis and Points for Cross-Validation}

The independent analysis reproduces the coherent story of Chapter~4 with three important wrinkles. For Research Question~1, self-reported familiarity and recent exposure are low (means near $1.9$; $22\%$ exposure), yet endorsement of specific capability descriptions is near-universal: a gap best described as low awareness coexisting with high capability-endorsement rather than as evidence of awareness as such. For Research Questions~2 and~3, perceived value and anticipated benefit are high and precisely estimated across the applicability, IT, and security blocks, with the caveat that the two investment-willingness items are at ceiling and contribute conviction but not discrimination. The comparative analyses are dominated by a very large Information Technology--versus--Information Assurance gap favoring Information Technology, and by a non-monotonic experience pattern in which mid-career respondents are the most reserved.

The following are the points where analytic choices most affect the conclusions, and therefore where two good-faith analyses of the same data are most likely to diverge; each is addressed in Chapter~4 or its limitations.
\begin{enumerate}[leftmargin=1.4em,itemsep=2pt]
  \item \textbf{Field operationalization.} The instrument yields four field categories, not two. Whether Engineering and Other are excluded, pooled, or analyzed as a four-group Kruskal--Wallis changes the comparative section; the analysis honors the proposal's IT-vs-IA contrast but reports all four groups in the focused breakouts.
  \item \textbf{Ceiling items.} Questions 4.5 and 5.7 are near-invariant. Treating them as ordinary outcomes versus acknowledging the ceiling materially changes how Research Question~2 receptiveness is reported; the analysis characterizes them as conviction rather than discrimination.
  \item \textbf{Awareness composite composition.} Averaging Q1.1 with Q2.1--Q2.6 blends two constructs that move in opposite directions; Q1.1 is therefore reported separately alongside the composite so its low central tendency is not obscured.
  \item \textbf{Experience shape.} The omnibus effect is real but non-linear; a monotonic narrative would misrepresent it. The $11$--$15$ year dip and the instability of the $16+$ cell ($n=18$) are reported explicitly.
  \item \textbf{Wave (non-response) bias.} The early-versus-late comparison shows no meaningful difference on any composite (all $|d|\le0.04$, $p>.5$), so response timing gives no evidence of self-selection; the upper-bound caveat therefore rests on the recruitment frame rather than on a wave signal, and Chapter~5 treats it as such.
\end{enumerate}

A fully scripted pipeline produced every statistic, figure, and table above from the response file, so each value is reproducible and directly checkable against an independent analysis of the same data. Taken together, the methodology, the worked examples, and this executable code and output leave the analysis of Chapter~4 fully traceable from raw responses to reported result.

